\section{Øvning}
\subsection{Guide til sangøvning for ikke-trænede sangere}
Revysangere, der ikke har sunget undefor revyregi øver typisk ved at sætte original-indspilningen af sangen på, og synge med til den.
Det er uhensigtsmæssigt af flere grunde:
Original-tonearten vil typisk være en anden end revy-tonearten, så de får ikke øvet "muscle memory" der er relevant når de når til scenen.
På originalindspilningen er der altid en lead-sanger der synger melodien, så det lyder altid godt. På den måde bliver de ikke konfronteret med hvordan de egentlig lyder, og hvad der skal arbejdes med.
Indspilningen kører også hele tiden videre, uanset om det lykkedes at synge det netop afspillede stykke eller ej. På den måde bliver fokus fjernet fra de stykker der har brug for øvning, og spredt ud over hele nummeret.
Ideelt burde sangere øve i den toneart sangen spilles i til revy, uden anden sang til og kun minimal instrumentering. De burde øve enkelte stykker specifikt og ikke gå videre i nummeret, indtil de er færdige med at arbejde med det gældende stykke; maksimum ét værs ad gangen.
Det er nemt at tage denne øvepraksis for givet hvis man er musiker, men det er ikke åbenlyst for folk der ikke har modtaget musikundervisning, så det er værd at sige, særligt til sine utrænede sangere.
Konkret kan man hjælpe med at implementere god øve-praksis ved at lave øvefiler til dem.
En øvefil er blot et lydklip med akkorder/akkompagnement til sangen i den relevante toneart, som sangerne kan øve sig på.
Lad der være en mængde vokal/melodi der er akkurat tilstrækkelig til at sangeren kan finde ind på sin stemme, men ikke mere.
Øveviler kan både laves som en simpel telefon-optagelse af akkompagnementet, eller ved at bruge track separation à la Moises og en DAW til at lave et backing track i den rigtige toneart med passende mængde vokal.